\subsection{Original source code}
\label{section:OriginalImplementation}
\textcite{rakuschek2025visual} offline implementation was written in Python and focuses on two core functions. \\
A time delay embedding function takes an 1D input signal (as a NumPy array) and the parameters $d$ (embedding dimension), $\tau$ (time delay), and $s$ (stride or window step size). By default tau and stride is set to one. The function uses a multidimensional array to store the individual windows of length $d$. The tau enforces the step size between the elements in each window, a larger value spreads the window further apart in time.
The stride determines how many samples to move forward before creating the next window, a larger stride produces fewer output points.
As shown in the supplemental material \textcite{Supplemental_document}, larger values for $\tau$ and s result in a thinner point cloud while the overall circular structure remains. \\
The plotting function takes an subplot, the multidimensional array resulting from the time delay embedding function and an title, which is per default "TDE", as parameters.
For the principal component analysis, the scikit-learn python library was used, the function essentially reduces the given time delay embedding output from its $d$ to two dimensions. In the next step, the scores which are the distances of each point from the origin are calculated, these euclidean distances from the origin are normalized to the range [0, 1] and the points are plotted with colors mapped to their normalized distances.